\SourceLeaf{233}%
rem huius argumenti de figuris generalissime conceptis expositionem ad aliam occasionem nobis reservare debemus.

\Article{7}
Investigemus iam formulam ad exprimendam mensuram curvaturae pro quovis puncto superficiei curvae. Denotante \(\mathrm{d}\sigma\) aream elementi huius superficiei, \(Z\mathrm{d}\sigma\) erit area projectionis huius elementi in planum coordinatarum \(x,y\); et perinde, si \(\mathrm{d}\Sigma\) est area elementi respondentis in superficie sphaerica, erit \(Z\mathrm{d}\Sigma\) area projectionis ad idem planum: signum positivum vel negativum ipsius \(Z\) vero indicabit situm projectionis similem vel oppositum situi elementi projecti: manifesto itaque illae projectiones eandem rationem quoad quantitatem, simulque eandem relationem quoad situm, inter se tenent, ut elementa ipsa. Consideremus iam elementum triangulare in superficie curva, supponamusque coordinatas trium punctorum, quae formant ipsius projectionem, esse
\[
\begin{array}{ll}
x,&y\\
x+\mathrm{d}x,&y+\mathrm{d}y\\
x+\delta x,&y+\delta y
\end{array}
\]
Duplex area huius trianguli exprimetur per formulam
\[
\mathrm{d}x\mathord{.}\delta y-\mathrm{d}y\mathord{.}\delta x
\]
et quidem in forma positiva vel negativa, prout situs lateris a puncto primo ad tertium respectu lateris a puncto primo ad secundum similis vel oppositus est situi axis coordinatarum \(y\) respectu axis coordinatarum \(x\).

Perinde si coordinatae trium punctorum, quae formant projectionem elementi respondentis in superficie sphaerica, a centro sphaerae inchoatae, sunt
\[
\begin{array}{ll}
X,&Y\\
X+\mathrm{d}X,&Y+\mathrm{d}Y\\
X+\delta X,&Y+\delta Y
\end{array}
\]
duplex area huius projectionis exprimetur per
\[
\mathrm{d}X\mathord{.}\delta Y-\mathrm{d}Y\mathord{.}\delta X
\]
de cuius expressionis signo eadem valent quae supra. Quocirca mensura curva%
